% AUTORIDAD DE FUENTE: Exposé IX, §§14.4--14.5, líneas 4726--4741;
% índices de escaneo 530--531, folios 519--520. Las capas francesas D/C
% no presentan diferencias materiales. Tres errores locales inequívocos
% se conservan reversiblemente mediante \sisi y se han comunicado al
% custodio raíz.

\par\medskip
\noindent
\underline{Conjetura 14.4.} Con las hipótesis que acabamos de precisar
en 14.3.5, existe un homomorfismo

\begin{equation}\tag{14.4.1}
 i:\underline M_U\longrightarrow G_U
\end{equation}

\noindent
tal que, para todo número primo $\ell$, la extensión (14.3.5.6) sea
isomorfa a la extensión deducida de $i$ como se explicó para (14.2.2).
Más precisamente, existe \sisi{una y una solo manera}{una y una sola
manera} de asociar a toda situación del tipo considerado en 14.3.5 un
homomorfismo (14.4.1) que satisfaga la condición precedente, de tal modo
que la formación de $i$ sea compatible con todo cambio de base
$S'\longrightarrow S$ asociado a un homomorfismo continuo
$V\longrightarrow V'$ de anillos noetherianos ádicos.

\medskip
\noindent
14.4.2. Parece que la existencia de $i$, de manera compatible con el
cambio de base, se deduce de la teoría analítica rígida de M. Raynaud,
habida cuenta de que la desarrolla en un marco más general que el
enunciado en 14.1, trabajando, como aquí, sobre anillos ádicos
noetherianos generales. Además, para demostrar la existencia y la
unicidad afirmadas en 14.4, nos reducimos inmediatamente a considerar
solo aquellos $S$ para los que $S_o$ es un esquema de tipo finito sobre
$\mathbb Z$. En efecto, $V$ es el límite inductivo filtrante de
subanillos $V_i$ de tipo finito sobre $\mathbb Z$; por el teorema de
Nagata (EGA IV 7.8), estos pueden suponerse normales y, para $i$
suficientemente grande, la situación $(A,A',W)$ proviene de una
situación análoga sobre $V_i$, o bien sobre el completado separado de
$V_i$ para la topología definida por
$m_i=\underline m V_i$. Esto da la reducción afirmada.

%%%%%%%%%%  scan idx 531  |  folio 520  %%%%%%%%%%

\medskip
\noindent
14.4.3. Obsérvese ahora que $S-S_o$ es un esquema de Jacobson
(EGA IV 10.5.8). Así, para ver que dos soluciones $i,i'$ del problema,
para una situación dada $(A,A',W)$, son iguales, basta demostrar que
estos homomorfismos son iguales en todo punto \underline{cerrado}
$\eta$ de $U$. Sea $S'$ la normalización de la clausura de $\eta$ en
$S$. Entonces $S'$ es un rasgo completo, y puede aplicarse el resultado
de unicidad de $\sisi{4.3}{14.3.4}$ a la imagen inversa sobre
$S'$ de la situación sobre $S$, pues el cuerpo residual de
$S'$ es \underline{finito}. Una vez establecida así la unicidad, esta
implica la compatibilidad de la formación de $i$ con el cambio de base,
a reserva de la existencia.

\medskip
\noindent
\underline{Observaciones 14.5.} Quedaría por formular con la generalidad
deseada, y por establecer, una compatibilidad entre el homomorfismo $i$
introducido en 14.4 por un método analítico rígido y el homomorfismo
análogo (13.5.13) de la teoría analítica compleja---o, mejor aún,
¡encontrar una generalización común de ambos!
